\chapter{Notación}
\label{ap:notacion}

Referencia rápida de la notación usada en todo el libro. Para
convenciones tipográficas (cajas, fuente monoespaciada, énfasis), ver
\emph{Convenciones} en la sección frontal.

\section{Grafos y vértices}

\begin{center}
\renewcommand{\arraystretch}{1.3}
\begin{longtable}{p{3.0cm}p{2.0cm}p{8.5cm}}
\toprule
\textbf{Símbolo} & \textbf{Tipo} & \textbf{Significado} \\
\midrule
\(G = (\V, \E)\)              & grafo            & grafo dirigido base; \(\V\) vértices, \(\E\) aristas. \\
\(G = (\V, \E, t)\)           & grafo temporal   & cada arista \(e \in \E\) lleva un \texttt{timestamp} \(t(e) \in \T\). \\
\(\V\)                        & conjunto         & vértices del grafo. \(|\V| = n\). \\
\(\E\)                        & multiconjunto    & aristas dirigidas; en multigrafos, varias \((u, v, t)\) pueden coexistir. \(|\E| = m\). \\
\(u, v, w, x\)                & vértices         & elementos genéricos de \(\V\). \\
\(e = (u, v, t)\)             & arista temporal  & del vértice origen \(u\) al destino \(v\), con \emph{timestamp} \(t\). \\
\(\T\)                        & dominio temporal & en el motor: \texttt{i64} segundos epoch (Unix time). \\
\(\Lc_V\), \(\Lc_E\)          & alfabeto         & etiquetas de vértices (\texttt{EntityType}) y aristas (\texttt{RelationType}). \(\Lc = \Lc_V \cup \Lc_E\). \\
\(\mathrm{type}(v)\)          & función          & \(\V \to \Lc_V\); el \texttt{EntityType} del vértice \(v\). \\
\(\mathrm{type}(e)\)          & función          & \(\E \to \Lc_E\); el \texttt{RelationType} de la arista \(e\). \\
\bottomrule
\end{longtable}
\end{center}

\section{Cardinalidades, grados y vecindarios}

\begin{center}
\renewcommand{\arraystretch}{1.3}
\begin{longtable}{p{3.0cm}p{2.5cm}p{8.0cm}}
\toprule
\textbf{Símbolo} & \textbf{Tipo} & \textbf{Significado} \\
\midrule
\(n = |\V|\)                  & escalar          & cantidad de vértices. \\
\(m = |\E|\)                  & escalar          & cantidad de aristas. \\
\(\deg(v)\)                   & escalar          & grado total de \(v\) (entradas + salidas). \\
\(\indeg(v)\)                 & escalar          & grado de entrada: \(|\{u : (u, v, \cdot) \in \E\}|\). \\
\(\outdeg(v)\)                & escalar          & grado de salida: \(|\{w : (v, w, \cdot) \in \E\}|\). \\
\(N(v)\)                      & subconjunto de \(\V\) & vecinos de \(v\) (en cualquier dirección). \\
\(N^{+}(v)\)                  & subconjunto de \(\V\) & vecinos salientes: \(\{w : (v, w, \cdot) \in \E\}\). \\
\(N^{-}(v)\)                  & subconjunto de \(\V\) & vecinos entrantes. \\
\(\bar{d}\)                   & escalar          & grado promedio: \(\bar{d} = m / n\). \\
\(\Delta\)                    & escalar          & grado máximo: \(\max_{v \in \V} \deg(v)\). \\
\bottomrule
\end{longtable}
\end{center}

\section{Hipótesis, joins, y enumeración}

\begin{center}
\renewcommand{\arraystretch}{1.3}
\begin{longtable}{p{3.0cm}p{2.5cm}p{8.0cm}}
\toprule
\textbf{Símbolo} & \textbf{Tipo} & \textbf{Significado} \\
\midrule
\(Q\), \(H\)                  & hipótesis        & consulta \emph{subgraph matching}; en GraphHunter, \texttt{Hypothesis}. \\
\(k\)                         & escalar          & longitud de la hipótesis: número de pasos / aristas. \\
\(\sigma\)                    & asignación       & mapeo de variables de la hipótesis a vértices del grafo. \\
\(R, R_i\)                    & relación         & multiconjunto de tuplas; en el motor, una \emph{bag} \(B_i\). \\
\(R \bowtie S\)               & relación         & \emph{join} natural sobre las columnas compartidas. \\
\(R \ltimes S\)               & relación         & \emph{semi-join}: filas de \(R\) cuya \emph{join key} aparece en \(S\). \\
\(\sigma_{P}(R)\)             & relación         & selección: filas de \(R\) que satisfacen el predicado \(P\). \\
\(\pi_{X}(R)\)                & relación         & proyección a las columnas \(X\). \\
\(\rhostar(Q)\)               & escalar          & \emph{fractional edge cover number} de \(Q\) (capítulo~\ref{cap:wcoj}). \\
\(\AGM(Q)\)                   & escalar          & cota \emph{Atserias-Grohe-Marx} sobre el tamaño de la salida: \(m^{\rhostar(Q)}\). \\
\bottomrule
\end{longtable}
\end{center}

\section{Índices y ventanas temporales}

\begin{center}
\renewcommand{\arraystretch}{1.3}
\begin{longtable}{p{3.4cm}p{2.5cm}p{7.6cm}}
\toprule
\textbf{Símbolo} & \textbf{Tipo} & \textbf{Significado} \\
\midrule
\(\NLF(v, L)\)                                & escalar             & vecinos salientes de \(v\) con etiqueta \(L \in \Lc_V\). \\
\(\bits[v]\)                                  & \texttt{u16}        & bitmap de tipos alcanzables desde \(v\) en \(\leq \texttt{MAX\_DEPTH}\) saltos. \\
\(\mathrm{min\_first}[v][L]\)                 & \texttt{i64}        & \emph{timestamp} mínimo del primer \emph{edge} de algún camino \(v \rightsquigarrow L\). \\
\(\mathrm{max\_first}[v][L]\)                 & \texttt{i64}        & análogo para el máximo. \\
\(\mathrm{rank}[v]\)                          & \texttt{u32}        & índice del vértice \(v\) en la fila \emph{sparse} de \texttt{min/max\_first}. \\
\([t_{\mathrm{lo}}, t_{\mathrm{hi}})\)        & intervalo en \(\T\) & ventana temporal de la consulta (cláusula \texttt{WITHIN}). \\
\(\Delta t = t_{\mathrm{hi}} - t_{\mathrm{lo}}\) & escalar          & ancho de la ventana. \\
\texttt{MAX\_DEPTH}                           & constante           & profundidad máxima de \emph{khop\_reach}: 4 en GraphHunter. \\
\texttt{NLF\_TYPE\_COUNT}                     & constante           & cantidad de \texttt{EntityType}s con \emph{tag} \([0, 9)\): \(|\Lc_V|\). \\
\bottomrule
\end{longtable}
\end{center}

\section{Tipos del motor (Rust) que aparecen recurrentemente}

\begin{center}
\renewcommand{\arraystretch}{1.3}
\begin{longtable}{p{4.0cm}p{9.0cm}}
\toprule
\textbf{Tipo} & \textbf{Significado} \\
\midrule
\texttt{Entity}              & nodo del grafo (público, humano): \texttt{id: String}, \texttt{entity\_type: EntityType}, \texttt{metadata}. \\
\texttt{Relation}            & arista del grafo (público): \texttt{source\_id}, \texttt{dest\_id}, \texttt{rel\_type}, \texttt{timestamp}, \texttt{metadata}. \\
\texttt{CompactRelation}     & arista compacta (interna): 32 bytes, \texttt{\#[repr(C)]}. Capítulo~\ref{cap:compact-relation}. \\
\texttt{StrId}               & \texttt{u32} compacto producido por el \emph{interner}. \\
\texttt{StringInterner}      & tabla bidireccional \texttt{String \(\leftrightarrow\) StrId}. \\
\texttt{Hypothesis}          & AST de una consulta DSL: \texttt{steps}, \texttt{shape}, \texttt{window\_seconds}, etc. \\
\texttt{HypothesisShape}     & enum \texttt{Chain | Graph}; lo lee el \emph{planner}. \\
\texttt{HuntResult}          & \texttt{Vec<Arc<str>>}; un \emph{path} resuelto al alfabeto humano. \\
\texttt{NlfTable}            & estructura del índice NLF; \emph{lazy build} en \texttt{OnceLock}. \\
\texttt{KHopReach}           & estructura del índice \emph{k-hop} con extensión temporal. \\
\texttt{IndexBacking<T>}     & \emph{wrapper} \emph{heap-resident} \(\leftrightarrow\) \emph{mmap-spilled} para tablas grandes. \\
\texttt{InMemoryStreamingBackend} & edge store \emph{degree-adaptive} per-vertex con auto-spill incremental. \\
\texttt{DispatchPlan}        & enum del \emph{planner}: \texttt{DfsMatch | K4Lftj | Yannakakis}. \\
\bottomrule
\end{longtable}
\end{center}

\section{Convenciones del pseudocódigo}

Reproducidas acá para referencia rápida (la versión narrativa está en
\emph{Convenciones}, sección frontal):

\begin{itemize}
    \item Asignación: \(x \leftarrow y\). Igualdad: \(x = y\). Acceso a campo: \(v.\mathrm{campo}\).
    \item \(\leftarrow\) y \(\Leftarrow\) ambos significan asignación; \(\Leftarrow\) sólo en derivaciones de variables agregadas.
    \item \texttt{\textbf{for/while/if/else/end}}: bloques delimitados por indentación + cierre explícito.
    \item \texttt{\textbf{return}}: sale del procedimiento; \texttt{\textbf{emit}}: agrega un resultado a la salida sin retornar.
    \item Conjuntos: \(\{x : P(x)\}\), \(A \cup B\), \(A \cap B\), \(A \setminus B\), \(\emptyset\).
    \item \emph{Bitwise}: \(a \mid b\), \(a \mathrel{|}\!= b\), \(a \,\&\, b\). Detalles cuando aparezcan.
\end{itemize}

\cleardoublepage
